% ======================================================================
%  TIER 2 KARTI: Offline vs Online Eval
% ======================================================================

\card{Offline vs Online Eval}{yayından önce ölç, yayından sonra öğren}
\cl{Problem} Golden dataset'te iyi skor almak yeterli mi? Gerçek kullanıcılar
farklı davranır. İki ayrı soru vardır: ``yayına hazır mı?'' ve ``gerçekte işe
yarıyor mu?'' --- ve bunlar farklı ölçümler ister.
\cl{Tanım} \textbf{Offline eval}: sabit golden dataset'te, yayından \emph{önce}
(hızlı, ucuz, regression yakalar). \textbf{Online eval}: gerçek trafikte, yayından
\emph{sonra} (A/B test, kullanıcı geri bildirimi, gerçek metrikler).
\cl{Model} Offline = simülatörde uçuş provası; online = gerçek yolcularla uçuş.
İkisi de gerekir; biri diğerinin yerini tutmaz.
\cl{Akış} Her değişiklik önce offline gate'ten geçer → sonra canary/A-B ile online
denenir. Online sinyaller (kabul oranı, düzeltme, terk) → yeni golden örnekler (flywheel).
\cl{``Yeterince iyi'' nedir} Kalite çıtası bir \emph{ürün kararıdır} --- \%100 değil,
hatanın maliyeti ile kullanışlılığın dengesi. Bunu PM tanımlar.
\cl{Tuzak} Yalnız offline'a güvenmek (gerçek dünya farklıdır); ya da yalnız online
(yavaş, riskli --- regression'ı canlıda öğrenirsin).
\cl{İlişkili} \term{Eval} · \term{Canary} · \term{Shadow Testing} · \term{Observability} · \term{Veri \& Gizlilik} (flywheel) \,{\scriptsize\color{softgray}(İlke \#13, \#18)}
